\phantomsection
\chapter{Процес виконання}
\section{Формування набору тестів}

Для оцінювання якості пошуку сформовано тестовий набір tests/eval-set.jsonl обсягом 15 пар <<запит -- релевантний чанк>>.

\textbf{Принцип відбору чанків і формулювання запитів:}
\begin{dashlist}[wide]
    \item Чанки відбиралися вручну з обох частин корпусу (data/chunks/wiki\_chunks.jsonl та data/chunks/quest\_chunks.jsonl) із різних предметних розділів і лінійок квестів (Tier05Steam, Tier1LV, MassProcessing, мультиблоки GregTech, бджільництво/дерева, інструменти Tinkers Construct та GT Tools) за умови наявності змістовного опису.
    \item До набору увійшли 7 чанків із книги квестів і 8 чанків зі статей вікі.
    \item Запити сформульовано англійською мовою (відповідно до мови індексованого корпусу) за двома типовими сценаріями використання системи:
    \begin{enumerate}[wide]
        \item пошук опису конкретного квесту або механіки (7 запитів -- наприклад, переваги та вимоги Spinmatron-2737, альтернативи переходу на LV Alloy Smelter, отримання саджанців у Tree Breeding);
        \item пошук унікального правила або пояснення, що трапляється лише в одному місці (8 запитів -- наприклад, заповнення вихідного люка Bacterial Vat на 50\%, заміна боросилікатного скла в Steam Fuser, ремонт інструмента Tinkers Construct, керування таблицею луту Dungeonchest.xml).
    \end{enumerate}
\end{dashlist}

\textbf{Чому обраний набір є показовим:}
\begin{dashlist}[wide]
    \item Поєднує два структурно різні типи текстів (енциклопедичні статті вікі та стислі квести зі збагаченими назвами предметів), що перевіряє роботу єдиного векторного індексу одразу по всьому корпусу.
    \item Містить запити на конкретні числові й технічні деталі, а не загальні описи, тому правильний чанк має бути знайдений точно серед тисяч інших (зокрема при збігу загальних слів на кшталт Alloy Smelter чи Tinkers Construct).
    \item Включає унікальні ігрові терміни (Spinmatron-2737, Dungeonchest.xml, Ball of Moss), що дозволяє перевірити стійкість ембединг-моделі до вузькоспеціалізованої лексики та коректність розрахунку метрик Recall@5, Recall@10 і MRR@10.
\end{dashlist}

\section{Обрання моделі для sentence-transformer ембедингів}

Для векторизації чанків через бібліотеку SentenceTransformers спочатку було протестовано багатомовну модель \texttt{BAAI/bge-m3}, яка підтримує кілька мов (зокрема українську й англійську). Проте під час тестування було виявлено, що вона є занадто великою за обсягом і споживанням ресурсів, при цьому на наявному наборі запитів не дає кращого результату пошуку, ніж компактніша англомовна модель.

З огляду на те, що весь підготовлений корпус даних і тестові запити сформовано англійською мовою, для подальшої роботи було обрано модель \texttt{BAAI/bge-base-en-v1.5}, яка забезпечує високу якість пошуку англомовних технічних текстів при значно менших витратах пам'яті та вищій швидкості індексування.

\section{Побудова векторного індексу}

Для побудови та обслуговування векторного індексу було розгорнуто сервер Qdrant у Docker-контейнері (із налаштуванням параметрів HNSW-індексу та payload-індексів у загальній конфігурації config/pipeline.toml).

Процес індексування реалізовано в модулях src/index/sentece\_indexig.py, src/index/vectorization.py та src/index\_main.py у такій послідовності:
\begin{dashlist}[wide]
    \item використано стандартні механізми бібліотеки SentenceTransformers для пакетної векторизації текстів (\texttt{encode\_texts}, \texttt{encode\_chunks\_jsonl}) з L2-нормалізацією векторів;
    \item векторизовано поле \texttt{text} усіх підготовлених чанків (wiki\_chunks.jsonl та quest\_chunks.jsonl);
    \item у класі QdrantVectorizer реалізовано формування структур \texttt{PointStruct} (із детермінованим UUID5-ідентифікатором, отриманим вектором і повним \texttt{payload} із метаінформацією чанку);
    \item реалізовано функції створення колекції, додавання payload-індексів та пакетного (batch) завантаження точок \texttt{upsert\_points} у Qdrant.
\end{dashlist}

\section{Пошук top-5}

Для перевірки роботи пошукового індексу через TUI-обгортку src/search.py для обраної моделі \texttt{BAAI/bge-base-en-v1.5} виконано серію тестових запитів у режимі top-5. За результатами перевірки всі потрібні релевантні чанки успішно знаходяться в перших 5 результатах видачі.

Нижче наведено використані запити для кожного тестового прогону:
\begin{dashlist}[wide]
    \item \textbf{Прогін 1 (див. рис.~\ref{fig:top5-test-1}):} запит <<How many parallels per voltage tier does the Thermic Heating Device provide and what is it useful for?>>;
    \item \textbf{Прогін 2 (див. рис.~\ref{fig:top5-test-2}):} запит <<Why does the mining level shown in the NEI Tool Materials tab differ from the extruder and smeltery tabs?>>;
    \item \textbf{Прогін 3 (див. рис.~\ref{fig:top5-test-3}):} запит <<Does a Ball of Moss repair a completely broken Tinkers Construct tool?>>.
\end{dashlist}

\begin{figure}[H]
    \centering
    \paddedimage[width=1\textwidth]{test_1-Thermic Heating Device.png}
    \caption{Тестовий прогін 1}
    \label{fig:top5-test-1}
\end{figure}

\begin{figure}[H]
    \centering
    \paddedimage[width=1\textwidth]{test_2-Tinkers Tools.png}
    \caption{Тестовий прогін 2}
    \label{fig:top5-test-2}
\end{figure}

\begin{figure}[H]
    \centering
    \paddedimage[width=1\textwidth]{test_3-Flint and Steel.png}
    \caption{Тестовий прогін 3}
    \label{fig:top5-test-3}
\end{figure}

\section{Тестування за набором}

Оцінювання якості пошуку проводилося скриптом tests/eval\_retrieval.py на підготовленому наборі tests/eval-set.jsonl (15 запитів) із розрахунком метрик Recall@5, Recall@10 та MRR@10. Для обох порівнюваних моделей використовувалися однакові параметри векторизації та пошуку в Qdrant (нормалізація векторів, метрика \texttt{cosine}, параметри HNSW-індексу, порожня інструкція запиту \texttt{query\_instruction} й обмеження \texttt{top\_k=10}), змінювалася лише сама модель ембедингів.

За результатами тестування (див. рис.~\ref{fig:eval-bge-v15} та рис.~\ref{fig:eval-bge-v3}) обидві моделі показали однакові підсумкові показники повноти (Recall@5 = 93.33\%, Recall@10 = 100.00\%), проте промах у топ-5 (\texttt{R@5 MISS}) припав на різні запити:
\begin{dashlist}[wide]
    \item \textbf{eval-10 для BAAI/bge-base-en-v1.5:} запит стосується оптимального рівня заповнення вихідного люка в Bacterial Vat (\textit{<<What is the optimal output hatch fill percentage for maximum growth and productivity in the Bacterial Vat?>>}). У цільовому чанку (\texttt{Bacterial Vat > Overview}, 1228 символів) перші близько 280 символів займають метадані інфобокса включно з довгим GitHub-посиланням, далі йде опис бактерій, Biolab і радіації, а ключова фраза про \textit{50\% capacity} міститься лише в самому кінці тексту. Через зміщення фокусу на початок чанку та наявність інших чанків про Bacterial Vat цей фрагмент опинився за межами топ-5 (але в межах топ-10).
    \item \textbf{eval-05 для BAAI/bge-m3:} запит перевіряє, чи ремонтує Ball of Moss повністю зламаний інструмент Tinkers Construct (\textit{<<Does a Ball of Moss repair a completely broken Tinkers Construct tool?>>}). Цільовий чанк походить із квесту \texttt{Tier05Steam > Flint and Steel}, де основний текст і списки предметів стосуються крафту кресала та подорожі в Nether, а потрібне правило згадане лише як побічна примітка з двох речень. Багатомовна модель \texttt{BAAI/bge-m3} у чистому dense-пошуку вище проранжувала профільні статті про ремонт Tinkers Tools, тоді як спеціалізована англомовна \texttt{BAAI/bge-base-en-v1.5} точніше виділила конкретне формулювання з примітки.
\end{dashlist}

\begin{figure}[H]
    \centering
    \paddedimage[width=1\textwidth]{BGE-v1.5-test.png}
    \caption{Тестовий прогін BAAI/bge-base-en-v1.5}
    \label{fig:eval-bge-v15}
\end{figure}

\begin{figure}[H]
    \centering
    \paddedimage[width=1\textwidth]{BGE-v3-test.png}
    \caption{Тестовий прогін BAAI/bge-m3}
    \label{fig:eval-bge-v3}
\end{figure}